SafeEdit-CRE围绕一个实际的调控基因组学问题而设计：当所选序列被第二个模型检验时，如何用少量替换改变一条天然顺式调控元件，同时保留细胞类型选择性信号？现有的序列-活性工作流提供了强大的预测器和日益成熟的设计工具，但候选排序往往与用作优化目标的同一个预测器绑定~\citep{Whalen2022,deBoer2024,Rafi2025}。SafeEdit-CRE将这些角色分开，并使用跨模型选择来决定哪些最小编辑值得推进。

与迭代贪婪优化最清晰的分野出现在600亲本评估中。SafeEdit-CRE将跨模型特异性裕度增益提高了0.055个标准化单位（95\% CI 0.040--0.071），尽管贪婪编辑已经在97.1\%的条件下产生了正裕度增益。因此，这种差异是在大多合理编辑中的排序效应：评审和不确定性标准偏向那些增益更大、且在不同模型间保持一致的变异体。这一效应在K562和HepG2的5--20核苷酸预算下最为显著，因为此时有更多编辑路径可用，且主模型偏好可能出现分化。

仅靠搜索广度并不能解释这一现象。计算量匹配的主束搜索的裕度增益为0.881，与SafeEdit-CRE的0.877高度吻合（差值$-0.004$；95\% CI $-0.011$--$0.004$）。然而SafeEdit-CRE返回的变异具有更低的跨模型不确定性（0.314对比0.335）和更小的靶标得分偏移。换言之，该框架保留了预测的细胞类型特异性目标，同时选择了候选图景中更不分散的部分。当合成和报告基因能力有限、实验预算更适合花在紧凑、可审计的面板上时，这种权衡尤其有价值。

九准则基准在更苛刻的设计设定下给出了同样的结论。SafeEdit-CRE的设计通过率为60.6\%，而贪婪优化为38.8\%，相差21.8个百分点（95\% CI 18.3--25.3）。完整重搜索消融揭示了原因：天然性感知重排序对不确定性降低的贡献最大，而序列域预过滤器在生成过程中强制执行GC和同源多聚体可行性。移除评审器精确重现了主束结果，表明跨模型评审是改变返回序列的组件，而不仅仅是在选择后对其做注释。

因此，输出是一个候选文库，而不是单一模型最优解。74条Tier A序列通过了全部九项准则，满足靶标和裕度要求，不带合成风险标记，且位于或接近主--评审帕累托前沿。它们覆盖三种靶细胞类型和四种编辑预算，提供了一个紧凑、表征充分的首轮面板，用于报告基因筛选和下游调控语法解析；1,141条Tier B序列则保留了广度，可用于更大规模的后续MPRA。候选记录保留了亲本序列、编辑后序列、编辑位置、模型特异性增益、不确定性、链一致性、天然性、碱基组成、同源多聚体状态和编辑计数。补充表~S10中的三个代表性示例凸显了具体的序列语法假说：局部替换在保留更广泛CRE骨架的同时改变细胞类型选择性，可通过报告基因实验直接检验。

若干设计选择使得比较具有可解释性。亲本选择是标签盲的，开发和评估亲本通过标识符和序列哈希分开，随机重复取平均而非按结果选择，置信区间对整个亲本进行重采样。冻结后评估器排除在候选生成之外，而计算量匹配的束搜索保持了搜索宽度和扩展成本不变。完整重搜索消融进一步确保了组件效应反映的是不同的返回序列，而非对一个候选表的追溯过滤。

最直接的生物学产出是一套紧凑、可审计的CRE变异，用于报告基因筛选。本研究的结果确立了计算排序及其跨模型行为；报告基因实验现在可以检验预测的裕度和不确定性排序是否与实测活性一致。由于候选保留了天然CRE骨架和明确的编辑预算，该面板也适用于将局部序列变化映射到细胞类型选择性调控输出的后续实验。除了报告基因筛选之外，最小编辑框架还具有直接的实际意义：Tier A面板可以指导合成基因电路中细胞类型特异性调控元件的设计，帮助优先选择疾病队列中功能性非编码变异进行实验验证，并为利用碱基编辑或先导编辑理性改造内源性CRE提供模板。每个候选都带有记录在案的编辑预算、链一致性检查和天然性评分，降低了选择那些合成失败或产生无法解释结果的序列的实验风险。

SafeEdit-CRE从而提供了一条从序列-活性预测到实验可检验的最小编辑的可重复路径。跨模型选择保留了主要的细胞类型特异性目标，相对于计算量匹配的单模型搜索降低了预测离散度，并产出了一个实用的Tier A面板，而非一长串不透明的无约束高分序列。同样的选择原则可以推广到其他调控预测模型和细胞情境——在这些地方，候选可靠性与名义最优值同样重要。
